\section{Universal Enclosure and Transfer Calculus}
\label{sec:universal-calculus}

This section collects the geometry shared by the transfer and direct-witness
arguments.  The same equilateral enclosure functional governs the local
admissible set of Strategy~2 and the forced witness set of Strategy~4.  The
same radical governs the signed center, all selected $T_+$ arcs, and the
standard local frontiers.  Strategy~2 uses one decorated $g$-family: a hat
imposes the nonsupercritical cap, $\vee$ takes the complementary
incoming-reach view, and extra subscripts or superscripts record center
intervals or lower relaxations.

\subsection{The equilateral enclosure gauge}

Let $\mathsf R$ denote counterclockwise rotation through $2\pi/3$.  For a
nonempty compact set $K\subset\mathbb R^2$, put
\[
 h_K(n)=\max_{x\in K}\langle x,n\rangle
\]
and
\begin{equation}
 \Lambda(K)=\frac2{\sqrt3}
 \min_{\lVert n\rVert=1}
 \sum_{j=0}^2h_K(\mathsf R^jn).
 \label{eq:universal-lambda}
\end{equation}

\begin{proposition}[Equilateral enclosure gauge]
\label{prop:universal-enclosure-gauge}
The number $\Lambda(K)$ is the least side length of a closed equilateral
triangle containing $K$.  It is monotone under inclusion, translation
invariant, and positively homogeneous.  If $K$ is a finite polygonal hull, a
minimizing normal may be chosen perpendicular to one of its hull edges.
\end{proposition}

\begin{proof}
For a fixed unit normal $n$, the three supporting half-planes with outward
normals $n,\mathsf Rn,\mathsf R^2n$ form an equilateral triangle.  Its side is
$2/\sqrt3$ times the sum of the three support numbers, which proves
\eqref{eq:universal-lambda}.  Monotonicity and homogeneity are immediate, while
translation invariance follows from
\[
 n+\mathsf Rn+\mathsf R^2n=0.
\]
On an open normal-direction cell with three fixed support points, the support
sum has the form $A\cos\theta+B\sin\theta$ and has second derivative equal to
its negative.  In the nondegenerate case an interior critical point is a
maximum.  Hence a minimum occurs at a cell boundary, where one support line
contains a hull edge.
\end{proof}

For the local demand hull
\[
 K(a,b,c)=\conv\{0,au,bv,c(u+v)\},
\]
where $u$ and $v$ are the two unit boundary directions meeting at $120^\circ$,
the admissible set is
\begin{equation}
 \mathcal A
 =\{(a,b,c)\in[0,1]^3:\Lambda(K(a,b,c))\le1\}.
 \label{eq:universal-admissible-gauge}
\end{equation}
The exact support cells and radial envelope are proved in
Appendix~\ref{sec:exact-local-demand-calculus}.

\subsection{One equilateral radical}

For $0\le x\le1$, define
\begin{equation}
 \omega(x)=\sqrt{1-x+x^2},
 \qquad
 \sigma(x)=1-\omega(x).
 \label{eq:universal-radical}
\end{equation}
Then
\begin{equation}
 \omega(1-x)=\omega(x),
 \qquad
 \sigma(x)=\frac{x(1-x)}{1+\omega(x)},
 \label{eq:universal-radical-identities}
\end{equation}
and direct differentiation gives
\begin{equation}
 \sigma''(x)=-\frac3{4\omega(x)^3}<0.
 \label{eq:universal-sigma-concavity}
\end{equation}
Thus $\sigma$ is strictly concave.  In the signed center calculus,
\[
 E=\omega(R),\qquad \eta=\sigma(R).
\]

For $0<x<1$, put $z=\sigma(x)/x$.  Eliminating the square root gives
\begin{equation}
 x=\frac{1-2z}{1-z^2},
 \qquad
 \sigma(x)=\frac{z(1-2z)}{1-z^2},
 \qquad 0<z<\frac12.
 \label{eq:universal-rational-radical}
\end{equation}
This is the common rational parameter used in the selected high-sheet
calculations.

\subsection{One \(g\)-family}

For $0\le x,c\le1$, define the historical defect-coordinate map
\begin{equation}
 g_c(x)
 =
 \max\{y\in[0,1]:(1-x,y,c)\in\mathcal A\},
 \qquad
 \widehat g_c(x)=\min\{g_c(x),x\}.
 \label{eq:reader-raw-capped-maps}
\end{equation}
Thus $x$ is the incoming boundary defect, $g_c(x)$ is the exact maximal
outgoing reach, and the hat imposes the nonsupercritical diagonal cap.

For any map $f:[0,1]\to[0,1]$, put
\begin{equation}
 f^\vee(a)=1-f(1-a).
 \label{eq:reader-complement-dual}
\end{equation}
Hence
\begin{equation}
 \widehat g_c^\vee(a)
 =
 1-\widehat g_c(1-a)
 \ge a.
 \label{eq:reader-capped-raw-relation}
\end{equation}

The exact appendix aliases are
\begin{equation}
 B_c(a)=g_c(1-a),
 \qquad
 F_c(a)=\widehat g_c(1-a),
 \qquad
 G_c(a)=\widehat g_c^\vee(a).
 \label{eq:reader-technical-aliases}
\end{equation}
They are retained only where they shorten the contact-cell algebra.

\begin{lemma}[Raw and capped actual-row transfer]
\label{lem:reader-raw-capped-transfer}
Suppose an actual row has incoming reach at least $a$, radial reach at least
$c$, and outgoing reach $B$.  If no center interval intervenes on the next
edge, then
\[
 B\le g_c(1-a),
 \qquad
 A_{\rm next}\ge g_c^\vee(a).
\]
If the row is nonsupercritical, then
\[
 B\le\widehat g_c(1-a),
 \qquad
 A_{\rm next}\ge\widehat g_c^\vee(a)\ge a.
\]
\end{lemma}

\begin{proof}
The same triangle realizes $(a,B,c)$, so the definition of $g_c$ gives
$B\le g_c(1-a)$.  A center-free boundary handoff gives
\[
 A_{\rm next}\ge1-B\ge1-g_c(1-a)=g_c^\vee(a).
\]
Under nonsupercriticality one also has $B\le1-a$, so
$B\le\min\{g_c(1-a),1-a\}=\widehat g_c(1-a)$.  Taking complements proves the
last assertion.
\end{proof}

\begin{proposition}[Zero-radial distinction]
\label{prop:reader-zero-radial-g}
For $0<x<1$,
\[
 g_0(x)>x,
 \qquad
 \widehat g_0(x)=x,
\]
and therefore
\[
 g_0^\vee(a)<a,
 \qquad
 \widehat g_0^\vee(a)=a
 \quad(0<a<1).
\]
\end{proposition}

\begin{proof}
The exact diameter formula is
\[
 B_0(a)=\frac{-a+\sqrt{4-3a^2}}2.
\]
Consequently
\[
 g_0(x)=B_0(1-x)
 =\frac{x-1+\sqrt{1+6x-3x^2}}2.
\]
Since
\[
 1+6x-3x^2-(1+x)^2=4x(1-x)>0,
\]
one has $g_0(x)>x$.  The hatted identity follows from its definition, and the
dual statements follow by complementation.
\end{proof}

\subsection{Center-assisted transfers and the free envelope}

Let $J\subseteq[0,1]$ be empty or a closed interval.  For $0\le p\le1$, put
\[
 e_J(p)=\sup\{x\in[0,1]:[0,x]\subseteq[0,p]\cup J\},
 \qquad
 \mathcal R_J(p)=1-e_J(p).
\]
If $J=[L,U]$, then
\begin{equation}
 \mathcal R_{[L,U]}(p)=
 \begin{cases}
  1-p,&p<L,\\
  1-\max\{p,U\},&p\ge L,
 \end{cases}
 \qquad
 \mathcal R_\varnothing(p)=1-p.
 \label{eq:reader-residual-formula}
\end{equation}
The map is nonincreasing in $p$ and under enlargement of $J$.

\begin{lemma}[Generalized boundary handoff]
\label{lem:reader-generalized-handoff}
Suppose a unit boundary edge is covered by a left trace contained in $[0,p]$,
a center trace $J$, and a right trace contained in $[1-q,1]$.  Then
\[
 q\ge\mathcal R_J(p).
\]
For an actual row as in Lemma~\ref{lem:reader-raw-capped-transfer}, define
\[
 g_{c,J}^\vee(a)
 =
 \mathcal R_J(g_c(1-a)).
\]
Then
\begin{equation}
 A_{\rm next}
 \ge \mathcal R_J(B)
 \ge g_{c,J}^\vee(a).
 \label{eq:reader-generalized-raw-transfer}
\end{equation}
If the row is nonsupercritical, define
\[
 \widehat g_{c,J}^\vee(a)
 =
 \mathcal R_J(\widehat g_c(1-a)).
\]
Then
\begin{equation}
 A_{\rm next}
 \ge \widehat g_{c,J}^\vee(a).
 \label{eq:reader-generalized-transfer}
\end{equation}
For $J=\varnothing$, these are $g_c^\vee$ and
$\widehat g_c^\vee$.
\end{lemma}

\begin{proof}
The component of $[0,p]\cup J$ containing $0$ is $[0,e_J(p)]$.  If
$1-q>e_J(p)$, the intervening open interval is uncovered.  Hence
$q\ge1-e_J(p)$.  The remaining assertions follow from the outgoing bounds in
Lemma~\ref{lem:reader-raw-capped-transfer} and the fact that $\mathcal R_J$ is
nonincreasing.
\end{proof}

For $0\le c<1/2$, define the single free strict-supercritical envelope
\begin{equation}
 g_c^{\rm sc}
 =
 \sup_{\{x:g_c(x)>x\}}g_c(x)
 =
 \frac{c+\sqrt{c^2-8c+4}}2.
 \label{eq:reader-free-supercritical-envelope}
\end{equation}
The supremum is not attained.  Therefore an actual strict-supercritical row
with radial demand at least $c$ satisfies
\begin{equation}
 B<g_c^{\rm sc},
 \qquad
 A_{\rm next}>1-g_c^{\rm sc}.
 \label{eq:reader-free-raw-envelope}
\end{equation}
The former notation was
$B_{\rm sc}(c)=g_c^{\rm sc}$ and
$A_{\rm sc}(c)=1-g_c^{\rm sc}$.

\subsection{Relaxed composition and path budgets}

For maps listed in geometric row order, write
\begin{equation}
 [\Phi_1\mid\cdots\mid\Phi_r](x)
 =
 (\Phi_r\circ\cdots\circ\Phi_1)(x).
 \label{eq:reader-chain-notation}
\end{equation}
The leftmost slot acts first.  Write $\mathrm I(x)=x$ and $\mathrm I^k$ for
$k$ consecutive identity slots.

\begin{lemma}[Relaxed composition]
\label{lem:reader-relaxed-composition}
Suppose actual demands satisfy
\[
 x_j\ge\Phi_j(x_{j-1})
 \qquad(1\le j\le r),
\]
where every $\Phi_j$ is nondecreasing.  If
$\underline\Phi_j\le\Phi_j$, then
\[
 x_r\ge
 [\underline\Phi_1\mid\cdots\mid\underline\Phi_r](x_0).
\]
\end{lemma}

\begin{proof}
Set $y_0=x_0$ and
$y_j=\underline\Phi_j(y_{j-1})$.  If $x_{j-1}\ge y_{j-1}$, then
\[
 x_j
 \ge\Phi_j(x_{j-1})
 \ge\Phi_j(y_{j-1})
 \ge\underline\Phi_j(y_{j-1})
 =y_j.
\]
Induction proves the claim.
\end{proof}

\begin{lemma}[Boundary-path budget]
\label{lem:reader-boundary-path}
Let $T_1,\ldots,T_k$ be roles on a chain of $k+1$ consecutive unit boundary
edges.  If the two external endpoint roles contribute at most $u$ and $v$ on
the first and last edges, then coverage forces
\begin{equation}
 \sum_{j=1}^k(A_j+B_j)\ge k+1-u-v.
 \label{eq:reader-path-budget}
\end{equation}
Consequently, if all $k$ internal rows are nonsupercritical, the chain is
impossible whenever $u+v<1$.
\end{lemma}

\begin{proof}
The first and last edges give $A_1\ge1-u$ and $B_k\ge1-v$.  Each middle edge
gives $B_j+A_{j+1}\ge1$.  Adding these $k+1$ inequalities proves
\eqref{eq:reader-path-budget}.
\end{proof}

\subsection{Low-root and order-transition parameters}

For the high-radial maps used later, define
\begin{equation}
 e(d)=\frac{1-d}{2}
 \left(1-\sqrt{4(1-d)^2-3}\right)
 \qquad
 \left(0<d<1-\frac{\sqrt3}{2}\right).
 \label{eq:signed-low-root-definition}
\end{equation}
For the low-radial order transition, put
\begin{equation}
 h(c)=\frac{2c}{1+\sqrt{16c^2-3}}
 \qquad\left(\frac12<c\le\frac{\sqrt3}{2}\right).
 \label{eq:signed-order-transition}
\end{equation}
The exact contact-cell calculation in
Appendix~\ref{sec:exact-local-demand-calculus} proves that $e(d)$ is the low
selected root for radial demand $1-d$ and that $h(c)$ is the transition
between the two selected $T$ orders.

\subsection{The universal selected-\(T_+\) curve}

\begin{proposition}[Universal selected-$T_+$ curve]
\label{prop:reader-universal-tplus}
Fix $0<d<1/2$, put $c=1-d$, and suppose an input $p$ belongs to a selected
$T_+$ arc of $\widehat g_c^\vee$.  Write
\[
 q=\widehat g_c^\vee(p),
 \qquad \nu=q-p,
 \qquad x=\frac{q-d}{1-d}.
\]
Then
\begin{equation}
 \nu=\sigma(x)=1-\sqrt{1-x+x^2},
 \label{eq:reader-psi}
\end{equation}
and
\begin{equation}
 q=d+(1-d)x,
 \qquad
 p=d+(1-d)x-\sigma(x).
 \label{eq:reader-tplus-normal}
\end{equation}
Consequently $p\mapsto\widehat g_{1-d}^\vee(p)$ is increasing and strictly
concave on every selected $T_+$ arc.
\end{proposition}

\begin{proof}
The selected quadratic is
\[
 (1-q)(q-d)=(1-d)^2\nu(2-\nu).
\]
Substitution of $q=d+(1-d)x$ gives
\[
 x(1-x)=\nu(2-\nu).
\]
The selected component uses the smaller root, namely $\nu=\sigma(x)$.  For
fixed $d$, the function
\[
 p_d(x)=d+(1-d)x-\sigma(x)
\]
is increasing and strictly convex by \eqref{eq:universal-sigma-concavity}.
Its inverse is increasing and strictly concave, and $q=d+(1-d)x$ is affine in
that inverse.
\end{proof}

With the low root in \eqref{eq:signed-low-root-definition}, the high-radial chord
is
\begin{equation}
 \widehat g_{1-d}^\vee(p)
 \ge p+\frac{d}{e(d)-d}(p-d).
 \label{eq:reader-high-chord}
\end{equation}
If $t=h(1-d)$ is the low-radial order transition, the corresponding chord is
\begin{equation}
 \widehat g_{1-d}^\vee(p)
 \ge p+\frac{1-2t}{t-d}(p-d).
 \label{eq:reader-low-chord}
\end{equation}

To avoid a new function letter, for every coefficient $\lambda$ proved valid
on a selected arc write
\begin{equation}
 \widehat g_{1-d}^{\vee,\lambda}(p)
 :=
 p+\lambda(p-d)
 \le
 \widehat g_{1-d}^\vee(p).
 \label{eq:reader-affine-relaxation}
\end{equation}
The superscript $\lambda$ denotes a lower relaxation, not a second exact map.

The rational form \eqref{eq:universal-rational-radical} gives, in the
historical notation $x=1-\beta$,
\begin{equation}
 \beta=\frac{z(2-z)}{1-z^2},
 \qquad
 m_\beta=\frac{1-z+z^2}{1-z^2}.
 \label{eq:reader-beta-rational}
\end{equation}

\subsection{Threshold routing}

For $0<d<1-\sqrt3/2$, the exact high-demand threshold is
\begin{equation}
 z>e(d)
 \quad\Longrightarrow\quad
 \widehat g_{1-d}^\vee(z)\ge1-e(d).
 \label{eq:reader-high-threshold}
\end{equation}
Define the decorated lower threshold transfer
\begin{equation}
 \widehat g_{1-d}^{\vee,\rm th}(z)
 =
 \begin{cases}
  z,&z\le e(d),\\
  1-e(d),&z>e(d).
 \end{cases}
 \qquad
 \widehat g_{1-d}^{\vee,\rm th}
 \le
 \widehat g_{1-d}^\vee.
 \label{eq:reader-threshold-transfer}
\end{equation}

\begin{lemma}[Threshold routing]
\label{lem:reader-threshold-routing}
Let $z_j=\Phi_j(z_{j-1})$, where every $\Phi_j$ is nondecreasing and extensive.
If $\Phi_k=\widehat g_{1-d}^\vee$ and $z_{k-1}>e(d)$, then
\[
 z_j\ge1-e(d)\qquad(j\ge k).
\]
If a chain contains later occurrences of
$\widehat g_{1-d_1}^\vee$ and $\widehat g_{1-d_2}^\vee$ and
\[
 z_0>e(d_1),
 \qquad
 \min\{e(d_1),e(d_2)\}<Q,
\]
then one of those two rows forces every subsequent output above $1-Q$.
\end{lemma}

\begin{proof}
Equation~\eqref{eq:reader-high-threshold} gives the first triggered output,
and extensivity preserves it.  For the two-threshold statement, if
$e(d_1)<Q$, trigger at the first row.  Otherwise
\[
 z_0>e(d_1)\ge Q>e(d_2),
\]
so trigger at the second row.
\end{proof}

The lemma is used only after an actual-row induction such as
\eqref{eq:reader-generalized-raw-transfer} or
\eqref{eq:reader-generalized-transfer} has identified every formal iterate as
a lower bound for the corresponding actual row.
